\subsubsection{Tuilage temporel sur accélérateur}
Le tuilage réduit le nombre de lancements GPU en calculant plusieurs lignes par bloc. Pour une correspondance en \(k>0\), la position lue dans la ligne précédente est distante de \(j-k+1\) colonnes. Soit \(R\) le maximum de ce décalage et de la portée diagonale égale à un. Pour produire une sortie \([u,v]\) après \(h\) lignes, la sous-étape \(t\) calcule \([\max(0,u-(h-t)R),v]\). Pour un bloc complet de largeur \(b\), la somme des largeurs donne au plus \(hb+Rh(h-1)/2\) cases.
\schema{tuile}{Le halo d'une tuile de trois lignes se réduit de deux colonnes par sous-étape jusqu'aux quatre colonnes de sortie.}{fig:tuile}{src/cuda/cuda.cu}
L'algorithme~\ref{alg:tile} précise les lectures globales et les échanges internes que représente la figure~\ref{fig:tuile}.
\begin{algorithm}[!htbp]
\caption{Avancement de plusieurs lignes dans un bloc CUDA}\label{alg:tile}
\begin{algorithmic}[1]
\Require Previous global row, output interval \([u,v]\), tile height \(h\le8\), predecessor reach \(R\)
\Ensure Updated global scores in \([u,v]\) after \(h\) rows
\For{\(t=1,\ldots,h\)}
  \State Read predecessors from global memory at \(t=1\) and from shared memory thereafter
  \State Update \([\max(0,u-(h-t)R),v]\) with Algorithm~\ref{alg:nvpd}
  \State Synchronize the threads of this block
\EndFor
\State Write the final scores in \([u,v]\) to global memory
\end{algorithmic}
\end{algorithm}
Les 256 fils d'un bloc produisent 256 ou 1\,024 colonnes. La grille choisit la largeur supérieure si sa taille moyenne atteint au moins deux blocs par multiprocesseur. Le calcul par ligne s'applique quand \(R>64\), quand la paire contient moins de huit lignes ou pour le traceback.
\chart{gpu/variantes_gpu.png}{Sur le GH200, la largeur de tuile modifie la durée du calcul sur les premières paires synthétiques à deux pour cent d'éditions injectées.}{fig:tuiles}{bench/gpu/cuda/profile_gpu.cu}
La figure~\ref{fig:tuiles} compare des appels après échauffement. Trois appels chauds fournissent une moyenne par processus. La dispersion entre les trois processus porte sur ces moyennes. À un million de bases et 2\,\% d'erreur, la tuile fixe prend \(5{,}351\pm0{,}088\)\,s et le noyau par ligne \(5{,}178\pm0{,}012\)\,s dans la même allocation.

\synthese{À un million de bases dans la campagne adaptative, les temps moyens passent de \(5{,}222\pm0{,}188\) à \(3{,}711\pm0{,}099\)\,s. Le rapport de 1,41 inclut le recalcul des colonnes du halo.}

\subsubsection{Lancements et accès à la mémoire de l’accélérateur}
Le coût d'une ligne dépend des lancements et des adresses lues par un groupe de fils. Nsight Systems, le profileur chronologique de NVIDIA, enregistre les appels hôtes et les opérations GPU. Dans ce profil, les durées des appels API hôtes peuvent recouvrir celles des opérations GPU.
\chart{gpu/trace_cuda.png}{Sept captures Nsight Systems après chauffe donnent les durées par domaine, leur écart-type et leur intervalle de confiance sur la paire de 5\,000 et 4\,992 bases.}{fig:lancements}{bench/repeated/trace/gpu_trace.py}
La figure~\ref{fig:lancements} inclut le noyau d'initialisation. Pour 4\,992 lignes de récurrence, le noyau de hauteur huit demande \(1+\lceil4992/8\rceil=625\) lancements. Les sept processus profilés par version exécutent auparavant un alignement hors capture. Tous donnent \(d=99\). Le temps cumulé des noyaux passe de \(6{,}2648\pm0{,}0087\) à \(3{,}2960\pm0{,}0004\)\,ms. Les copies prennent \(8{,}00\pm0{,}79\) et \(8{,}39\pm0{,}77\)\,\(\mu\mathrm{s}\). Les appels API hôtes totalisent \(42{,}77\pm70{,}99\) et \(9{,}55\pm9{,}24\)\,ms. Dans l'ancienne version, un processus cumule 203,14\,ms d'API. L'intervalle de Student sur cette moyenne est très large et sa borne basse franchit zéro. L'intervalle exprime ici l'imprécision de l'estimation.
Pour un accès logique \(l\) d'un warp \(w\) de 32 fils, les adresses \(a\) des fils actifs se regroupent en secteurs de 32 octets. Le volume demandé par case se calcule par
\[
 B=\frac{32}{mn}\sum_{w,l}\bigl|\{\lfloor a/32\rfloor : a\in A_{wl}\}\bigr|
 \quad [\mathrm{bytes/cell}].
\]
Le modèle énumère les adresses de l'index, les lectures de scores et les écritures selon les positions des bases dans les séquences. Deux instructions qui accèdent au même secteur comptent séparément. L'accès aux compteurs L1 et L2 du GH200 est restreint sur cette allocation. Le trafic qui atteint la mémoire HBM reste indéterminé.
\chart{gpu/localite_gpu.png}{La capacité des secteurs demandés se partage entre les mots distincts de chaque instruction et les emplacements inutilisés. Le total est exprimé en octets par case.}{fig:secteurs}{bench/gpu/locality/gpu_sectors.cu}
Sur la paire de portée 30, la figure~\ref{fig:secteurs} donne 22,48 octets par case pour le noyau par ligne et 10,62 pour la tuile. Lorsque la portée atteint 169, le sélecteur reprend le noyau par ligne.

\synthese{Sur le rectangle de 32 lignes, la tuile réduit de 64,2\,\% le volume de secteurs demandés. Ce modèle porte sur les adresses émises et laisse ouverte la part des demandes finalement servies par la mémoire HBM.}

\subsubsection{Durée de vie des allocations sur accélérateur}
Une matrice de \(N\) séquences contient \(P=N(N-1)/2\) paires non ordonnées. L'allocation indépendante des deux tampons GPU requiert \(2P\) appels à \texttt{cudaMalloc}. Avec un espace de travail persistant, chaque tampon conserve son allocation jusqu'au prochain changement de taille. Si \(b_{it}\) désigne la taille positive du tampon \(t\) pour la paire \(i\), le nombre d'allocations physiques du lot vaut
\[
 A=\sum_{t=1}^2\left(1+\sum_{i=2}^{P}[b_{it}\ne b_{i-1,t}]\right).
\]
Cette expression suppose des entrées non vides et le traitement sériel des paires dans leur ordre de lecture.
\chart{gpu/allocations_cuda.png}{Pour chaque lot, le gain de temps est apparié sur trois blocs. Les deux nombres à gauche indiquent les appels à \texttt{cudaMalloc} avant et avec réutilisation.}{fig:alloc}{bench/gpu/allocation/allocator_measure.py}
Pour les 64 phages de la figure~\ref{fig:alloc}, la sonde compte 4\,032 puis 250 allocations physiques. Le pic reste de 150\,892 octets. Dans la sonde après chauffe, le temps passe de 13,346 à 8,527\,s, soit un rapport apparié de 1,563.

\synthese{Dans l’allocation distincte 838238, le temps du CLI complet sur ce lot passe de \(23{,}32\pm1{,}26\) à \(9{,}35\pm0{,}22\)\,s avec la réutilisation des tampons. Leur durée de vie réduit les allocations sans changer le pic mémoire du lot.}

\subsection{Alignement par fronts d’onde sur accélérateur}
Pour la distance à coût unitaire, eWFA-GPU d’Aguado-Puig et coll.~\cite{wfagpu,ewfacode} distribue les paires entre blocs et avance les diagonales d'un même front en parallèle. Son format de trace réserve deux bits à chacune des 64 opérations d'édition possibles. La boucle de cette version s'arrête avant le score 64. L'expérience retient les cas \(d<64\) parmi vingt paires définies avant le calcul. Pour chaque chemin brut, le programme rejoue les opérations jusqu'aux deux extrémités et vérifie
\[
 N_{\mathrm{sub}}+N_{\mathrm{ins}}+N_{\mathrm{del}}=d_{\rm edlib}.
\]
Les compteurs \(N_{\mathrm{sub}},N_{\mathrm{ins}},N_{\mathrm{del}}\) dénombrent respectivement substitutions, insertions et délétions dans le chemin.
\chart{gpu/front_gpu.png}{Sur les paires admissibles, les temps de processus comparent trois méthodes avec reconstruction du chemin sur les mêmes lots.}{fig:ewfa}{bench/gpu/ewfa/ewfa.sbatch}
La figure~\ref{fig:ewfa} porte sur cinq processus par lot. Pour cinq paires de 5\,000 bases faiblement divergentes, le processus eWFA-GPU prend \(297{,}21\pm0{,}58\)\,ms, dont \(2{,}931\pm0{,}034\)\,ms dans son chronomètre interne. NvPD CUDA prend \(379{,}22\pm5{,}47\)\,ms et edlib sur un cœur Grace \(6{,}24\pm0{,}57\)\,ms, du lancement du processus à sa fin. Sur les quatorze paires admissibles, le score concorde avec edlib et le rejeu du chemin brut consomme les deux chaînes à ce coût. Le formateur de sortie du programme amont produit quatorze CIGAR qui échouent au rejeu.

\synthese{Les quatorze chemins bruts admissibles retrouvent les scores attendus. Les cinq paires de 5\,000 bases à 2\,\% coûtent de 97 à 99 éditions et dépassent la capacité de la trace de cette version.}

\subsubsection{Répartition des fronts d’onde entre accélérateur et processeur}\label{sec:wfa-gpu-batch}
WFA-GPU, développé par Aguado-Puig et coll.~\cite{wfagpuaffine,wfagpucode}, distribue les alignements affines entre blocs CUDA et les diagonales d’un front entre fils. Soient \(x\) le coût d’une substitution, \(o\) celui de l’ouverture d’une lacune et \(e\) celui de son extension. Pour un chemin \(P\), le coût dépend de l’ensemble des lacunes \(G(P)\) et de leurs longueurs \(\ell_g\),
\[
 C(P)=xN_{\mathrm{sub}}(P)+\sum_{g\in G(P)}(o+e\ell_g),\qquad
 (x,o,e)=(1,0,1)\ \Longrightarrow\ \min_P C(P)=d(A,B).
\]
Le mode score utilise le front complet. Le GPU traite au plus \(512\) paires par lot. Lorsqu’un front atteint la limite de calcul ou qu’une chaîne contient \texttt{N}, WFA-GPU confie la paire à son sous-module WFA2~\cite{wfa2gpucode}. Le symbole \texttt{N} conserve sa valeur littérale.

La figure~\ref{fig:wfagpuaffine} compare cinq blocs après un bloc de chauffe sur GH200. Chaque temps couvre le lancement du processus, la lecture des chaînes et l’écriture des scores. À \(20\)~knt et \(10\,\%\) d’éditions injectées, le rapport des temps WFA-GPU/WFA2 passe de \(15{,}77\) pour cinq paires à \(0{,}762\) pour les \(128\) paires qui les contiennent. Sur ce dernier lot, WFA-GPU prend \(0{,}395\pm0{,}006\)~s, edlib \(0{,}250\pm0{,}001\)~s et NvPD CUDA \(2{,}524\pm0{,}028\)~s. NvPD parcourt le rectangle de chaque paire, tandis que la largeur du front dépend du score atteint.
\chart{gpu/lot_fronts_affines.png}{Rapport des temps de processus WFA-GPU et WFA2 sur un cœur Grace pour les lots synthétiques. Les points donnent les moyennes des rapports appariés, les traits épais l’écart-type et les traits fins l’intervalle de Student à 95\,\%. Le panneau b répartit les paires biologiques entre CPU et GPU. Les effectifs sont identiques dans les cinq blocs.}{fig:wfagpuaffine}{python/studies/wfa_gpu_batch.py}
\FloatBarrier
Sur les \(2\,016\) paires de \(\phi\)X174, les trois génomes contenant \texttt{N} participent à \(186\) paires. Le programme les reprend toutes sur CPU. Sur le chromosome~21, les \(768\) fenêtres homologues de \(16\,384\) bases de référence ont des distances comprises entre \(0\) et \(11\,107\). Porter la limite \(K\) de \(64\) à \(1\,024\) fait passer le nombre de reprises de \(152\) à \(10\). Le rapport apparié \(T_{1024}/T_{64}=0{,}958\) a pour intervalle à 95\,\% \([0{,}899\,;\,1{,}017]\). Le réglage change aussi la taille de bloc que le programme choisit pour le GPU.

Pour la matrice SARS-CoV-2, le pilote réserve le GPU aux paires dont chaque chaîne compte moins de \(32\,768\) bases et appartient à \(\{A,C,G,T\}^{*}\). Les six génomes admis fournissent \(\binom{6}{2}=15\) paires sur \(2\,016\). Quatre travailleurs CPU utilisant WFA2-lib 2.3.6~\cite{wfa2code} calculent les \(2\,001\) autres paires avec les symboles IUPAC littéraux, pendant que le GPU calcule les quinze scores. Les témoins CPU utilisent les mêmes quatre travailleurs. Du FASTA au TSV, le calcul hybride prend \(6{,}844\pm0{,}015\)~s, WFA2 \(6{,}874\pm0{,}030\)~s et edlib \(2{,}628\pm0{,}014\)~s. L’intervalle du rapport hybride/WFA2 vaut \([0{,}9908\,;\,1{,}0004]\).

À coût unitaire, le front de score \(s\) admet au plus \(2s+1\) diagonales, avant troncature aux extrémités des chaînes. Pour une paire de distance \(d\), le nombre cumulé de positions possibles vaut donc \(H(d)=\sum_{s=0}^{d}(2s+1)=(d+1)^2\). Les quinze paires admissibles sur GPU ne représentent que \(223\,071\) des \(26\,084\,927\,678\) positions de ce modèle. Celui-ci omet les extensions de correspondances. Dans la mesure, le lot GPU prend en moyenne \(0{,}498\)~s et le lot CPU le plus lent \(6{,}684\)~s.

\synthese{Sur le lot synthétique le plus divergent à \(128\) paires de \(20\)~knt, le rapport à WFA2 a pour intervalle à 95\,\% \([0{,}748\,;\,0{,}776]\). Sur les fenêtres chromosomiques, le meilleur temps moyen de WFA-GPU reste \(14{,}12\) fois celui de WFA2 sur un cœur Grace.}

\subsubsection{Portage de la récurrence par Kokkos}
Les espaces d'exécution et de mémoire de Kokkos appliquent l'algorithme~\ref{alg:nvpd} à la même représentation des scores sur EPYC et GH200. Deux \texttt{View} portent l'index et les lignes de DP. Dans le parcours par ligne, une même instance d'exécution ordonne un premier \texttt{parallel\_for} d'initialisation puis \(m\) appels de mise à jour. Une clôture précède la lecture du score. Le noyau natif OpenMP garde son équipe entre les lignes et le natif CUDA peut fusionner plusieurs lignes.
\chart{distributed/portage_kokkos.png}{Sur chaque architecture, une allocation fournit la moyenne de trois processus après échauffement. Les intervalles portent sur trois allocations.}{fig:kokkos}{bench/kokkos/kokkos_trial.cpp}
Les mesures du parcours par ligne dans la figure~\ref{fig:kokkos} incluent les allocations et les transferts de la paire. Sur huit cœurs EPYC, le calcul prend \(109{,}526\pm1{,}388\)\,ms avec Kokkos/OpenMP et \(27{,}953\pm0{,}038\)\,ms en OpenMP natif pour la référence de 20\,000 bases. Sur le même GH200, trois allocations comparent Kokkos~4.5.01 avec CUDA~12.6.3 et le noyau CUDA natif. L'écart-type porte sur les trois moyennes d'allocation, comme sur EPYC. À 5\,000 bases, le natif prend \(4{,}150\pm0{,}023\)\,ms et Kokkos \(11{,}131\pm0{,}792\)\,ms.

\synthese{Sur GH200, les mêmes entrées de 20\,000 bases prennent \(46{,}230\pm3{,}335\)\,ms avec Kokkos par ligne et \(17{,}179\pm0{,}123\)\,ms avec les tuiles natives. Les deux versions réservent 480\,024 octets sur le GPU.}

\subsubsection{Équipe persistante et fusion de lignes sur chromosomes}\label{sec:kokkoshumain}
Le parcours Kokkos de la figure~\ref{fig:kokkos} soumet une mise à jour par ligne. Sur l'EPYC, une équipe peut rester active pendant tout l'alignement. Chaque fil garde un intervalle contigu de colonnes. Une barrière rend les scores visibles avant de passer à la ligne suivante. Le parcours par ligne reste disponible lorsque le nombre de fils dépasse la taille d'équipe autorisée.

Sur GPU, le mode score peut fusionner deux lignes en un seul lancement. Soit \(U_x(P)\) la ligne obtenue en appliquant la récurrence NvPD à la lettre \(x\) et à la ligne précédente \(P\). Pour deux lettres successives de \(A\),
\[
 D_{i+1,\cdot}=U_{A_{i+1}}\!\left(U_{A_i}(D_{i-1,\cdot})\right),\qquad
 N_{\mathrm{updates}}=\left\lceil\frac{m}{2}\right\rceil .
\]
Chaque colonne finale recalcule au plus trois scores de la ligne intermédiaire, même si ces colonnes sont très éloignées. Une dernière ligne isolée conserve sa mise à jour simple. Le mode trace sur GPU stocke chaque ligne pour la remontée et garde le parcours par ligne.

\chart{distributed/kokkos_humain.png}{L'accélération est le rapport du temps natif au temps Kokkos sur la même cible. Le troisième panneau compare le RSS et les tableaux de travail des deux cibles.}{fig:kokkoshumain}{bench/kokkos/human/measure.py}
Les comparaisons portent sur deux régions de GRCh38 et CHM13. À 16\,384 bases de référence, la région colinéaire coûte 14 éditions. L'autre comprend une insertion de 9\,540 bases dans CHM13 et coûte 9\,568 éditions. Chacune des 66 configurations emploie cinq processus indépendants. Chaque processus calcule deux fois le score de la même paire et chronomètre le second appel. Dans les dix fenêtres, la portée de l'index des dernières occurrences varie de 86 à 1\,813 colonnes. Le noyau CUDA natif utilise donc son parcours par ligne. La fusion Kokkos recalcule ses prédécesseurs à leur adresse et accepte ces portées.

Dans le panneau a) de la figure~\ref{fig:kokkoshumain}, sur huit cœurs EPYC, la fenêtre colinéaire de 65\,536 bases prend \(1{,}0801\pm0{,}0050\)\,s avec Kokkos par ligne. L'équipe persistante ramène ce temps à \(0{,}6520\pm0{,}0004\)\,s. Le noyau OpenMP natif prend \(0{,}2841\pm0{,}0012\)\,s. Dans le panneau b), sur la fenêtre colinéaire de 524\,288 bases, CUDA natif prend \(1{,}990\pm0{,}035\)\,s. Kokkos par ligne donne \(2{,}520\pm0{,}034\)\,s et la fusion \(1{,}345\pm0{,}018\)\,s. Sur la fenêtre de même longueur avec insertion, la fusion donne \(1{,}376\pm0{,}014\)\,s.

Sur la fenêtre colinéaire de 65\,536 bases, une capture Nsight Systems compte 65\,537 lancements pour le parcours par ligne et 32\,769 avec la fusion, initialisation comprise. Les temps GPU cumulés de cette capture unique valent 95,83 et 64,06\,ms. Dans les deux cas, quatre transferts déplacent 1\,049\,300 octets.

Le panneau c) donne le pic RSS, la taille des tableaux hôtes et, sur GH200, les octets réservés sur le GPU pour les noyaux natif et Kokkos retenu. Sur la fenêtre de 65\,536 bases avec insertion, chacune des trois implémentations atteint un pic RSS de 180\,Mio dans ses cinq processus et réserve 1\,802\,376 octets sur le GPU. L'appel courant de Kokkos emploie l'équipe persistante sur l'hôte lorsque sa taille est admise, et les lignes fusionnées sur le GPU pour le score seul.

\FloatBarrier
\synthese{Sur les fenêtres humaines du GH200, fusionner deux lignes accélère Kokkos d’un facteur de 1,71 à 1,87. La portée des prédécesseurs impose au noyau CUDA natif de traiter ces mêmes fenêtres ligne par ligne.}

